% Generated draft: requires mathematical and source review.
\documentclass[11pt]{article}
\usepackage[margin=25mm]{geometry}
\usepackage{fontspec}
\setmainfont{Latin Modern Roman}
\usepackage{amsmath,amssymb,mathtools,unicode-math}
\setmathfont{Latin Modern Math}
\usepackage{enumitem,needspace,xurl,hyperref}
\setlength{\parindent}{0pt}
\setlength{\parskip}{5pt}
\setlength{\emergencystretch}{3em}
\newcommand{\N}{\mathbb{N}}
\newcommand{\Z}{\mathbb{Z}}
\newcommand{\Q}{\mathbb{Q}}
\newcommand{\R}{\mathbb{R}}
\newcommand{\C}{\mathbb{C}}
\newcommand{\F}{\mathbb{F}}
\newcommand{\A}{\mathbb{A}}
\newcommand{\PP}{\mathbb{P}}
\newcommand{\E}{\mathbb{E}}
\newcommand{\eps}{\varepsilon}
\newcommand{\dd}{\,\mathrm{d}}
\begin{document}
\section*{UnsolvedMath ID: 30002274}
\textbf{Source problem code:} catalogue entry
\subsection{Definitions and mathematical statement}
A numerical graph invariant is a numerical quantity assigned to a graph that is unchanged under graph isomorphism (equivalently, under relabeling of vertices). More formally, for a specified class of graphs \(\mathcal G\) and a specified numerical value space \(D\), an invariant is a function \(I\colon\mathcal G\to D\) satisfying
\[
G\cong H\ \Longrightarrow\ I(G)=I(H)
\qquad(G,H\in\mathcal G).
\]
Here \(\mathcal G\) and \(D\) are placeholders: the supplied record specifies neither. The question also permits numerical data without explicitly requiring isomorphism invariance.

\textbf{Source question.} What kinds of numerical data or invariants should be extracted from very large graphs, such as protein-interaction graphs, to describe useful structural information?

This is an open-ended methodological question. The source does not specify the graph model, a size threshold for ``very large,'' which structural features are relevant, or a criterion for usefulness. Protein-interaction graphs are given as an example, not as an exclusive domain. Consequently, the supplied information does not determine a quantified proposition, a particular invariant to construct, or a precise optimization or verification target. These qualifications remain unresolved.
\subsection{Short English statement}
The question asks which numerical summaries of very large graphs, including protein-interaction graphs, convey useful structural information. The intended graph class and meaning of usefulness are unspecified, so this remains a methodological research question requiring clarification before a precise theorem or construction task can be stated.
\subsection{Sources}
\begin{itemize}
\item Local registry: \texttt{data/unsolvedmath/problems.json}, permanent UnsolvedMath ID \(30002274\), title ``Structural Invariants for Very Large Graphs,'' source problem code \texttt{catalogue entry}; supplied collection: Oberwolfach Reports Open Problems.
\item Supplied report citation: \emph{Graph Theory} (2013), Oberwolfach Reports, pp.\ 67--128, DOI \texttt{10.4171/owr/2013/02}. Source URL: \url{https://doi.org/10.4171/owr/2013/02}. The record describes Oberwolfach Reports as documenting mathematical workshops and their open-problem discussions. No individual proposer is supplied.
\item Supplied literature source, labeled ``original Oberwolfach report'': \url{https://doi.org/10.4171/owr/2013/02}.
\item Literature assessment dated \texttt{2026-08-21}: the supplied curation labels the problem open and describes it as an open-ended methodological question with no definite numerical invariant or proposition to verify. This is supplied dated curation, not an independently verified literature assessment; time-sensitive status and cited claims require independent verification.
\item The supplied record labels the statement ``exact'' and reports that its mathematical content was checked against the official report, with notation, context, and question form normalized without changing the source-level claim. This verification claim is supplied provenance and has not been independently checked for this draft.
\end{itemize}
\end{document}
